\input prelude
% tex.ch's hash_extra (changes/web2c.ch): names hash modulo hash_prime=8501
% into hash_size=15000 slots, and a name whose slot is taken goes above
% eqtb_size. \hashx, \agzyz, \agzzx, \ahxyz, \ahxzx, \ahywz and \ahyxx all
% hash to 3228, so all but the first live in the extra region. Each path
% tex.ch changes for it is used: a new name (§260), its meaning and
% assignment tracing (§252), a local value restored at the end of a group
% (§283), \let and \def of such a name (§1215), \csname and \ifcsname.
\def\hashx{one}
\def\agzyz{two}
\def\agzzx{three}
\message{2007 [\hashx] [\agzyz] [\agzzx]}
\message{2007 \meaning\hashx}
\message{2007 \meaning\agzyz}
\message{2007 \meaning\agzzx}
{\def\agzyz{inner}\message{2007 [\agzyz]}\global\let\ahxyz\agzyz}
\message{2007 [\agzyz] [\ahxyz]}
\let\ahxzx\agzzx
\message{2007 \meaning\ahxzx}
\expandafter\def\csname ahywz\endcsname#1{<#1>}
\message{2007 [\ahywz{x}]}
\expandafter\ifx\csname ahyxx\endcsname\relax \message{2007 relax}\fi
\ifcsname ahywz\endcsname \message{2007 defined}\else \message{2007 undefined}\fi
\begingroup \let\hashx\undefined \message{2007 \meaning\hashx}\endgroup
\message{2007 \meaning\hashx}
\setbox0=\hbox{}
\lsshipbox0
\end
